\subsection{Foundations of Associative Mapping}

A dictionary is Python's built-in associative mapping object. It stores relationships between keys and values. The key is the lookup handle. The value is the object associated with that key.

A sequence answers a positional question:

\begin{quote}
What object is stored at index \texttt{2}?
\end{quote}

A dictionary answers an associative question:

\begin{quote}
What value is associated with this key?
\end{quote}

This changes the shape of access. A dictionary is not mainly a row of components. It is a mapping from a hashable key domain to an arbitrary value domain.

\begin{verbatim}
quotes = {
    "Homer": "D'oh!",
    "Jerry": "What's the deal with airline food?",
    "Arthur": "It is only a model.",
}

print(f"quotes: {quotes!r}")
print(f"type(quotes): {type(quotes)}")
print(f"quotes['Homer']: {quotes['Homer']!r}")
\end{verbatim}

Possible output:

\begin{verbatim}
quotes: {'Homer': "D'oh!", 'Jerry': "What's the deal with airline food?", 'Arthur': 'It is only a model.'}
type(quotes): <class 'dict'>
quotes['Homer']: "D'oh!"
\end{verbatim}

The expression \texttt{quotes['Homer']} does not mean ``the item at position \texttt{'Homer'}.'' It means: use the key \texttt{'Homer'} to retrieve the value associated with that key.

\subsubsection{Structural Syntax: The Key-Value Paradigm and Curly-Brace \texttt{\{\}} Literals}

A dictionary literal uses curly braces with key-value pairs inside. Each pair has the form:

\begin{verbatim}
key: value
\end{verbatim}

Several pairs are separated by commas.

\begin{verbatim}
person = {
    "name": "Homer",
    "job": "safety inspector",
    "age": 42,
}

print(f"person: {person!r}")
print(f"type(person): {type(person)}")
\end{verbatim}

Possible output:

\begin{verbatim}
person: {'name': 'Homer', 'job': 'safety inspector', 'age': 42}
type(person): <class 'dict'>
\end{verbatim}

The keys in this dictionary are \texttt{"name"}, \texttt{"job"}, and \texttt{"age"}. The values are \texttt{"Homer"}, \texttt{"safety inspector"}, and \texttt{42}.

Access is performed through the key.

\begin{verbatim}
person = {
    "name": "Homer",
    "job": "safety inspector",
    "age": 42,
}

print(f"name: {person['name']!r}")
print(f"job:  {person['job']!r}")
print(f"age:  {person['age']!r}")
\end{verbatim}

Possible output:

\begin{verbatim}
name: 'Homer'
job:  'safety inspector'
age:  42
\end{verbatim}

The dictionary object exposes methods that reflect mapping behavior.

\begin{verbatim}
person = {
    "name": "Homer",
    "job": "safety inspector",
    "age": 42,
}

print(f"type(person): {type(person)}")
print(f"__getitem__ in dir: {'__getitem__' in dir(person)}")
print(f"__setitem__ in dir: {'__setitem__' in dir(person)}")
print(f"__contains__ in dir: {'__contains__' in dir(person)}")
print(f"keys in dir: {'keys' in dir(person)}")
print(f"values in dir: {'values' in dir(person)}")
print(f"items in dir: {'items' in dir(person)}")
print(f"get in dir: {'get' in dir(person)}")
\end{verbatim}

Possible output:

\begin{verbatim}
type(person): <class 'dict'>
__getitem__ in dir: True
__setitem__ in dir: True
__contains__ in dir: True
keys in dir: True
values in dir: True
items in dir: True
get in dir: True
\end{verbatim}

The crucial dictionary operations are:

\begin{itemize}
    \item \texttt{\_\_getitem\_\_}: retrieves a value by key, as in \texttt{d[key]};
    \item \texttt{\_\_setitem\_\_}: assigns a value to a key, as in \texttt{d[key] = value};
    \item \texttt{\_\_contains\_\_}: checks key membership with \texttt{in};
    \item \texttt{keys}: returns a dynamic view of the dictionary keys;
    \item \texttt{values}: returns a dynamic view of the dictionary values;
    \item \texttt{items}: returns a dynamic view of key-value pairs;
    \item \texttt{get}: retrieves a value safely when the key may be absent.
\end{itemize}

Assigning to an existing key replaces the associated value. It does not create a second copy of the key.

\begin{verbatim}
person = {
    "name": "Homer",
    "job": "safety inspector",
}

print(f"before: {person!r}")

person["job"] = "donut inspector"

print(f"after:  {person!r}")
print(f"len(person): {len(person)}")
\end{verbatim}

Possible output:

\begin{verbatim}
before: {'name': 'Homer', 'job': 'safety inspector'}
after:  {'name': 'Homer', 'job': 'donut inspector'}
len(person): 2
\end{verbatim}

The dictionary still has two keys. The value associated with \texttt{"job"} changed.

Assigning to a new key inserts a new key-value pair.

\begin{verbatim}
person = {
    "name": "Homer",
    "job": "safety inspector",
}

person["quote"] = "D'oh!"

print(f"person: {person!r}")
print(f"len(person): {len(person)}")
\end{verbatim}

Possible output:

\begin{verbatim}
person: {'name': 'Homer', 'job': 'safety inspector', 'quote': "D'oh!"}
len(person): 3
\end{verbatim}

The syntax boundary between sets and dictionaries is important. A colon creates dictionary entries.

\begin{verbatim}
one_set = {"spam"}
one_dict = {"food": "spam"}
empty_dict = {}

print(f"one_set: {one_set!r}")
print(f"type(one_set): {type(one_set)}")
print()

print(f"one_dict: {one_dict!r}")
print(f"type(one_dict): {type(one_dict)}")
print()

print(f"empty_dict: {empty_dict!r}")
print(f"type(empty_dict): {type(empty_dict)}")
\end{verbatim}

Possible output:

\begin{verbatim}
one_set: {'spam'}
type(one_set): <class 'set'>

one_dict: {'food': 'spam'}
type(one_dict): <class 'dict'>

empty_dict: {}
type(empty_dict): <class 'dict'>
\end{verbatim}

The empty curly-brace literal \texttt{\{\}} creates an empty dictionary, not an empty set.

\subsubsection{Integrity Constraints: Why Dictionary Keys Must Be Hashable and Hash-Stable}

Dictionary lookup is hash-based. The key must be hashable because Python uses the key's hash value to locate the internal table area where the key-value pair belongs.

Immutable built-in objects such as strings, integers, and tuples of hashable elements can be used as dictionary keys.

\begin{verbatim}
data = {
    "name": "Arthur",
    42: "answer",
    ("x", "y"): "coordinate labels",
}

print(f"data: {data!r}")
print(f"type(data): {type(data)}")
print()

print(f"data['name']: {data['name']!r}")
print(f"data[42]: {data[42]!r}")
print(f"data[('x', 'y')]: {data[('x', 'y')]!r}")
\end{verbatim}

Possible output:

\begin{verbatim}
data: {'name': 'Arthur', 42: 'answer', ('x', 'y'): 'coordinate labels'}
type(data): <class 'dict'>

data['name']: 'Arthur'
data[42]: 'answer'
data[('x', 'y')]: 'coordinate labels'
\end{verbatim}

The keys can be inspected with \texttt{hash()}.

\begin{verbatim}
keys = [
    "name",
    42,
    ("x", "y"),
]

for key in keys:
    print(f"key={key!r:<12} type={type(key)} hash={hash(key)}")
\end{verbatim}

Possible output:

\begin{verbatim}
key='name'      type=<class 'str'> hash=5524778262942691467
key=42          type=<class 'int'> hash=42
key=('x', 'y')  type=<class 'tuple'> hash=7801337796950902376
\end{verbatim}

The exact hash values for strings and tuples may differ between Python runs. The important point is that each key can produce a hash value during the current execution.

Mutable built-in containers such as lists cannot be used as dictionary keys.

\begin{verbatim}
bad_key = ["name"]

print(f"bad_key: {bad_key!r}")
print(f"type(bad_key): {type(bad_key)}")
print(f"bad_key.__hash__: {bad_key.__hash__}")

try:
    data = {bad_key: "Homer"}
except TypeError as err:
    print(f"dictionary construction failed: {err}")
\end{verbatim}

Possible output:

\begin{verbatim}
bad_key: ['name']
type(bad_key): <class 'list'>
bad_key.__hash__: None
dictionary construction failed: unhashable type: 'list'
\end{verbatim}

A list is unhashable because its contents can change. If a list could be used as a key, the dictionary could place it according to one content state and later fail to find it after the list changed.

\begin{verbatim}
items = ["spam", "eggs"]

print(f"before: {items!r}")
items.append(42)
print(f"after:  {items!r}")
\end{verbatim}

Possible output:

\begin{verbatim}
before: ['spam', 'eggs']
after:  ['spam', 'eggs', 42]
\end{verbatim}

Dictionary keys must be hash-stable while stored. This means the data used by their hash and equality behavior must not change while the key is inside the dictionary.

A tuple can be used as a key only if all its elements are hashable.

\begin{verbatim}
safe_key = ("Homer", 42)
bad_key = ("Homer", [42])

print(f"safe_key: {safe_key!r}")
print(f"hash(safe_key): {hash(safe_key)}")
print()

print(f"bad_key: {bad_key!r}")

try:
    print(hash(bad_key))
except TypeError as err:
    print(f"hash failed: {err}")
\end{verbatim}

Possible output:

\begin{verbatim}
safe_key: ('Homer', 42)
hash(safe_key): -9029866530700280856

bad_key: ('Homer', [42])
hash failed: unhashable type: 'list'
\end{verbatim}

The outer tuple is immutable, but the inner list is not hashable. Since tuple equality and hashing depend on the contained elements, the tuple cannot be used as a safe key.

Two equal keys identify the same dictionary entry.

\begin{verbatim}
key1 = ("Jerry", "George")
key2 = ("Jerry", "George")

data = {}

data[key1] = "first version"
data[key2] = "second version"

print(f"key1 == key2: {key1 == key2}")
print(f"hash(key1) == hash(key2): {hash(key1) == hash(key2)}")
print(f"data: {data!r}")
print(f"len(data): {len(data)}")
\end{verbatim}

Possible output:

\begin{verbatim}
key1 == key2: True
hash(key1) == hash(key2): True
data: {('Jerry', 'George'): 'second version'}
len(data): 1
\end{verbatim}

The second assignment replaced the value associated with the existing equal key. It did not create a second independent mapping entry.

The key rule is the same rule used by sets:

\begin{quote}
Dictionary keys must be hashable, and their hash-relevant state must remain stable while they are stored.
\end{quote}

\subsubsection{Value Flexibility: Storing Arbitrary, Nested Data Types and Mutable Heap Objects}

Dictionary values are much less restricted than dictionary keys. A value does not need to be hashable because Python does not use the value to place the entry in the hash table. The key controls lookup. The value is the payload associated with that key.

Values can be strings, numbers, lists, dictionaries, sets, tuples, or other objects.

\begin{verbatim}
profile = {
    "name": "Lisa",
    "age": 8,
    "quotes": ["If anyone wants me, I'll be in my room."],
    "scores": {"math": 10, "music": 10},
    "tags": {"student", "saxophone"},
}

print(f"profile: {profile!r}")
print(f"type(profile): {type(profile)}")
print()

print(f"type(profile['name']): {type(profile['name'])}")
print(f"type(profile['age']): {type(profile['age'])}")
print(f"type(profile['quotes']): {type(profile['quotes'])}")
print(f"type(profile['scores']): {type(profile['scores'])}")
print(f"type(profile['tags']): {type(profile['tags'])}")
\end{verbatim}

Possible output:

\begin{verbatim}
profile: {'name': 'Lisa', 'age': 8, 'quotes': ["If anyone wants me, I'll be in my room."], 'scores': {'math': 10, 'music': 10}, 'tags': {'student', 'saxophone'}}
type(profile): <class 'dict'>

type(profile['name']): <class 'str'>
type(profile['age']): <class 'int'>
type(profile['quotes']): <class 'list'>
type(profile['scores']): <class 'dict'>
type(profile['tags']): <class 'set'>
\end{verbatim}

This is legal because only the keys must be hashable. The values may be mutable.

A dictionary can store a reference to a mutable object. If that object is mutated later, the dictionary still refers to the same object.

\begin{verbatim}
quotes = ["D'oh!", "Mmm... donuts"]

profile = {
    "name": "Homer",
    "quotes": quotes,
}

print(f"profile before: {profile!r}")
print(f"id(quotes): {id(quotes)}")
print(f"id(profile['quotes']): {id(profile['quotes'])}")

quotes.append("Why you little!")

print()
print(f"profile after:  {profile!r}")
print(f"id(quotes): {id(quotes)}")
print(f"id(profile['quotes']): {id(profile['quotes'])}")
\end{verbatim}

Possible output:

\begin{verbatim}
profile before: {'name': 'Homer', 'quotes': ["D'oh!", 'Mmm... donuts']}
id(quotes): 2199812427072
id(profile['quotes']): 2199812427072

profile after:  {'name': 'Homer', 'quotes': ["D'oh!", 'Mmm... donuts', 'Why you little!']}
id(quotes): 2199812427072
id(profile['quotes']): 2199812427072
\end{verbatim}

The dictionary did not copy the list automatically. It stored a reference to the same list object.

The list object can be inspected separately.

\begin{verbatim}
quotes = ["D'oh!", "Mmm... donuts"]

print(f"type(quotes): {type(quotes)}")
print(f"append in dir: {'append' in dir(quotes)}")
print(f"__getitem__ in dir: {'__getitem__' in dir(quotes)}")
print(f"__len__ in dir: {'__len__' in dir(quotes)}")
\end{verbatim}

Possible output:

\begin{verbatim}
type(quotes): <class 'list'>
append in dir: True
__getitem__ in dir: True
__len__ in dir: True
\end{verbatim}

The dictionary value is a normal list object. It keeps its own list behavior even though it is stored inside a dictionary.

Nested dictionaries are common because they allow named access at several levels.

\begin{verbatim}
characters = {
    "Homer": {
        "show": "The Simpsons",
        "quote": "D'oh!",
    },
    "Arthur": {
        "show": "Monty Python",
        "quote": "It is only a model.",
    },
}

print(f"type(characters): {type(characters)}")
print(f"type(characters['Homer']): {type(characters['Homer'])}")
print()

print(f"Homer quote: {characters['Homer']['quote']!r}")
print(f"Arthur show: {characters['Arthur']['show']!r}")
\end{verbatim}

Possible output:

\begin{verbatim}
type(characters): <class 'dict'>
type(characters['Homer']): <class 'dict'>

Homer quote: "D'oh!"
Arthur show: 'Monty Python'
\end{verbatim}

The expression \texttt{characters['Homer']['quote']} performs two dictionary lookups:

\begin{verbatim}
characters['Homer']       -> inner dictionary
inner_dictionary['quote'] -> final string value
\end{verbatim}

This is one of the main reasons dictionaries are central in Python. They allow a program to represent named structure without relying on numeric positions.

\subsubsection{Membership Testing: Why \texttt{x in d} Checks Keys, Not Values}

For dictionaries, the expression \texttt{x in d} checks whether \texttt{x} is a key in the dictionary. It does not check whether \texttt{x} is one of the values.

\begin{verbatim}
quotes = {
    "Homer": "D'oh!",
    "Jerry": "No soup for you!",
    "Arthur": "It is only a model.",
}

print(f"quotes: {quotes!r}")
print()

print(f"{'Homer' in quotes = }")
print(f"{'D\\'oh!' in quotes = }")
\end{verbatim}

Possible output:

\begin{verbatim}
quotes: {'Homer': "D'oh!", 'Jerry': 'No soup for you!', 'Arthur': 'It is only a model.'}

'Homer' in quotes = True
"D'oh!" in quotes = False
\end{verbatim}

The value \texttt{"D'oh!"} exists inside the dictionary, but it is not a key. Therefore, \texttt{"D'oh!" in quotes} is \texttt{False}.

This behavior follows from the dictionary's architecture. A dictionary is optimized to find values by key. The key is hashed. The value is not used for table placement.

The membership operation is equivalent to asking whether the key exists.

\begin{verbatim}
quotes = {
    "Homer": "D'oh!",
    "Jerry": "No soup for you!",
    "Arthur": "It is only a model.",
}

print(f"'Homer' in quotes:        {'Homer' in quotes}")
print(f"'Homer' in quotes.keys(): {'Homer' in quotes.keys()}")
\end{verbatim}

Possible output:

\begin{verbatim}
'Homer' in quotes:        True
'Homer' in quotes.keys(): True
\end{verbatim}

The first form is preferred because it is direct and idiomatic.

Checking values requires using \texttt{.values()} explicitly.

\begin{verbatim}
quotes = {
    "Homer": "D'oh!",
    "Jerry": "No soup for you!",
    "Arthur": "It is only a model.",
}

print(f"{'D\\'oh!' in quotes.values() = }")
print(f"{'Homer' in quotes.values() = }")
\end{verbatim}

Possible output:

\begin{verbatim}
"D'oh!" in quotes.values() = True
'Homer' in quotes.values() = False
\end{verbatim}

This is a different operation. Key lookup uses the dictionary hash table. Value membership must inspect values, because values are not the hash-table lookup domain.

The view objects can be inspected.

\begin{verbatim}
quotes = {
    "Homer": "D'oh!",
    "Jerry": "No soup for you!",
    "Arthur": "It is only a model.",
}

keys_view = quotes.keys()
values_view = quotes.values()
items_view = quotes.items()

print(f"type(keys_view): {type(keys_view)}")
print(f"type(values_view): {type(values_view)}")
print(f"type(items_view): {type(items_view)}")
print()

print(f"keys_view: {keys_view}")
print(f"values_view: {values_view}")
print(f"items_view: {items_view}")
\end{verbatim}

Possible output:

\begin{verbatim}
type(keys_view): <class 'dict_keys'>
type(values_view): <class 'dict_values'>
type(items_view): <class 'dict_items'>

keys_view: dict_keys(['Homer', 'Jerry', 'Arthur'])
values_view: dict_values(["D'oh!", 'No soup for you!', 'It is only a model.'])
items_view: dict_items([('Homer', "D'oh!"), ('Jerry', 'No soup for you!'), ('Arthur', 'It is only a model.')])
\end{verbatim}

The details of dictionary view objects will be developed later. For now, the important distinction is simple:

\begin{itemize}
    \item \texttt{x in d}: checks keys;
    \item \texttt{x in d.keys()}: also checks keys, but is usually unnecessary;
    \item \texttt{x in d.values()}: checks values;
    \item \texttt{x in d.items()}: checks key-value pair tuples.
\end{itemize}

A final example shows all three domains.

\begin{verbatim}
quote_map = {
    "Homer": "D'oh!",
    "Jerry": "No soup for you!",
}

print(f"{'Homer' in quote_map = }")
print(f"{'D\\'oh!' in quote_map = }")
print(f"{'D\\'oh!' in quote_map.values() = }")
print(f"{('Homer', 'D\\'oh!') in quote_map.items() = }")
\end{verbatim}

Possible output:

\begin{verbatim}
'Homer' in quote_map = True
"D'oh!" in quote_map = False
"D'oh!" in quote_map.values() = True
('Homer', "D'oh!") in quote_map.items() = True
\end{verbatim}

The dictionary has three visible logical domains: keys, values, and key-value pairs. Bare membership testing belongs to the key domain.